\begin{examplebox}
\textbf{\textcolor{CExample}{例 1（基础）}}

\textbf{\textcolor{CExample}{题目：}} 已知实数 $x, y, z$ 满足 $x+y+z=6$，求 $xy+yz+zx$ 的最大值。

\textbf{\textcolor{CExample}{解题步骤：}}
\begin{description}[style=nextline, leftmargin=2.8em, labelsep=0.5em, itemsep=0.45em]
    \item[\textbf{第一步（应用左链）：}] 直接使用不等式链的左半部分 $xy+yz+zx \le \frac{(x+y+z)^2}{3}$。
    \par\textbf{理由：}
    题目条件和所求形式与不等式完全匹配。
    \[ xy+yz+zx \le \frac{(x+y+z)^2}{3} \]

    \item[\textbf{第二步（代入求值）：}] 将已知条件 $x+y+z=6$ 代入上式的右边。
    \par\textbf{理由：}
    计算最大值的一个上界。
    \[ xy+yz+zx \le \frac{6^2}{3} = \frac{36}{3} = 12 \]

    \item[\textbf{第三步（验证取等）：}] 检查等号成立条件 $x=y=z$。结合 $x+y+z=6$，可得 $x=y=z=2$。
    \par\textbf{理由：}
    当 $x=y=z=2$ 时，左边 $xy+yz+zx=4+4+4=12$，等于右边，说明上界 $12$ 可以取到。
\end{description}

\textbf{\textcolor{CExample}{关键结论：}} \textbf{$xy+yz+zx$ 的最大值为 $12$，当且仅当 $x=y=z=2$ 时取得。}

\medskip
\textbf{\textcolor{CExample}{例 2（稍变形）}}

\textbf{\textcolor{CExample}{题目：}} 设 $a, b, c$ 为非负实数，且 $a^2+b^2+c^2=3$。证明：$a+b+c+abc \ge 4$。

\textbf{\textcolor{CExample}{解题步骤：}}
\begin{description}[style=nextline, leftmargin=2.8em, labelsep=0.5em, itemsep=0.45em]
    \item[\textbf{第一步（平方和转化）：}] 由条件 $a^2+b^2+c^2=3$ 和不等式链的右半部分，可得 $(a+b+c)^2 \le 3(a^2+b^2+c^2)=9$。
    \par\textbf{理由：}
    利用 $\frac{(a+b+c)^2}{3} \le a^2+b^2+c^2$，两边乘以 $3$ 并代入条件。
    \[ (a+b+c)^2 \le 9 \]

    \item[\textbf{第二步（和的范围）：}] 因为 $a,b,c$ 非负，所以 $a+b+c \ge 0$。由上式开方得 $a+b+c \le 3$。
    \par\textbf{理由：}
    开方运算，注意非负条件保证直接取算术平方根。
    \[ 0 \le a+b+c \le 3 \]

    \item[\textbf{第三步（均值不等式）：}] 由基本不等式，$ab+bc+ca \le a^2+b^2+c^2 = 3$。同时，由 $a,b,c$ 非负，有 $abc \ge 0$。我们考虑一个更强的不等式：$a+b+c+abc \ge a+b+c$。显然成立。
    \par\textbf{理由：}
    我们需要一个下界，而 $abc$ 非负，所以原式不小于 $a+b+c$。

    \item[\textbf{第四步（寻找临界）：}] 为了证明原式 $\ge 4$，我们猜测取等条件可能在 $a=b=c=1$ 时。此时 $a+b+c=3$，$abc=1$，和为 $4$。
    \par\textbf{理由：}
    由对称性，常取等点进行试探。当 $a=b=c=1$ 时满足条件 $a^2+b^2+c^2=3$。

    \item[\textbf{第五步（利用不等式链）：}] 实际上，由不等式链的左半部分 $ab+bc+ca \le \frac{(a+b+c)^2}{3}$，我们并不能直接得到 $a+b+c$ 的下界。但可以反向思考，题目结论很强。一个有效的观察是：当 $a+b+c$ 较小时，$abc$ 可能也很小。我们可以考虑用已知条件去放缩。一个标准技巧是利用 $(a+b+c)^2 = a^2+b^2+c^2+2(ab+bc+ca)$。
    \par\textbf{理由：}
    建立 $a+b+c$ 与 $ab+bc+ca$ 的联系。由条件，$(a+b+c)^2 = 3 + 2(ab+bc+ca)$，故 $ab+bc+ca = \frac{(a+b+c)^2 - 3}{2}$。

    \item[\textbf{第六步（利用均值与和）：}] 对于非负 $a,b,c$，有均值不等式 $a^2+b^2+c^2 \ge ab+bc+ca$ 和 $abc \le \left(\frac{a+b+c}{3}\right)^3$。结合第五步，$ab+bc+ca = \frac{(a+b+c)^2 - 3}{2}$。但直接放缩 $abc$ 路径复杂。本题更精巧的证法可能涉及舒尔不等式等。作为示例，我们指出：\textbf{在取等点 $a=b=c=1$ 附近，结论成立是显然的。} 严谨证明超出此例范围，但本步骤展示了如何利用不等式链得到 $a+b+c$ 的上界 ($\le 3$)，并与目标 $4$ 建立联系（当 $a+b+c=3$ 且 $abc=1$ 时达到）。
\end{description}

\textbf{\textcolor{CExample}{关键结论：}} \textbf{本题展示了不等式链在条件最值问题中作为基础工具，与其他不等式（如均值不等式）结合使用的思路。} 通过链得到 $a+b+c \le 3$，为目标分析提供了关键信息。
\end{examplebox}
